% Status figure; diagnostic numbers, not the final certification table.
% Centralized slide overview. Regenerate with: python figures/gen_k1_centralized_overview.py
\begin{figure*}[t]
  \centering
  \includegraphics[width=0.98\textwidth]{figures/k1_centralized_overview.pdf}
  \caption{Centralized SpaceMiT K1 performance overview with one aligned row per model in every
  panel. Warm latency and within-row peer ratios compare Merlin against ExecuTorch for the INT8
  rows and against the same-graph XNNPACK GEMM swap for the FP32 VLA rows. Operator attribution
  and one-to-eight-core scaling retain the same seven-row layout and explicitly mark measurements
  that are not yet available. Cross-run, timeout-bound, and profiler-qualified results remain
  visibly distinguished rather than being presented as certified peers.}
  \label{fig:k1-centralized-overview}
\end{figure*}

% Status figure; diagnostic numbers, not the final certification table.
\begin{figure*}[t]
  \centering
  \includegraphics[width=0.95\textwidth]{figures/k1_executorch_status.pdf}
  \caption{SpaceMiT K1 performance status; lower latency is better. (a) Historical FP32 bitVLA
  experiment using the same Merlin graph/runtime, with only 15 FP32 matmuls routed to the named
  kernel backend. Points are three independent launches and horizontal marks are minima. Merlin's
  148.3 ms was 11.4\% lower latency than the XNNPACK-kernel hybrid in this restricted experiment;
  this was not a full ExecuTorch comparison. (b) Diagnostic one- and eight-core full-flow latency
  ratio $T_{\mathrm{ET}}/T_{\mathrm{Merlin}}$ for W8A8/qd8; values below 1 mean Merlin is slower.
  These are best valid diagnostic walls, not a paired uniform certification set. The smolVLA
  one-core triangle is an upper bound because Merlin did not complete within 120 s; its ExecuTorch
  reference is a seeded-random recomputation rather than Merlin's captured weights. smolVLA has no
  matched eight-core measurement. ResNet-50 passes the diagnostic absolute W8A8 bar but lacks an
  independent FP32 golden. (c)
  Complete-output TinyLlama diagnostic points at one and eight cores, grouped by compiler
  configuration rather than splicing the best point from different binaries into one curve. The
  MR4/NR32 model-vectorizer-enabled pair uses the exact same ELF and has three-launch median latency
  1244.6 ms at one core and 402.1 ms at eight cores, only 1.05$\times$/1.08$\times$ the fresh
  ExecuTorch warm slopes. Its same-session MR4/NR16 control is 1583.5/474.8 ms. The interleaved
  MR8/NR16 candidate regresses to 2175.7/458.6 ms because wider accumulator residency increases
  scalar operand-feed pressure. Both ExecuTorch points are fresh N=1/N=4 warm slopes. The validated Merlin configurations produced
  the reported 256k-output SHA and pass the shared accuracy bar. SHA agreement establishes repeatability, not
  reference correctness. Governor was performance at 1.6 GHz for the current INT8 measurements;
  core counts were matched and multicore affinity was bounded to the requested cores. No
  across-model INT8 performance win is currently certified. (d) Full-output, correctness-gated
  LSTMNetVIT scaling. Replacing four grouped-im2col paths with direct grouped convolution and
  suppressing OpenMP regions below 10k static work units, prequantizing all 17 structurally eligible
  constant contraction weights AOT, and folding 39 residual broadcast-to-add/mul chains reduces
  Merlin from 433.0 to a 130.9 ms launch median at one core and from 86.9 to 49.6 ms at eight cores.
  Three feature-on launch means from an off/on/on/off/off/on full-output A/B are shown per current
  Merlin point; each launch contains five timed iterations after two warmups. The distinct medians
  across all 15 raw iterations are 131.1/49.5 ms. Merlin remains 2.52$\times$ and 3.23$\times$ slower
  than ExecuTorch at one and eight cores, respectively.}
  \label{fig:k1-executorch-status}
\end{figure*}

% Current-best Merlin warm latency and per-op attribution for VLA workloads.
% Regenerate with: python figures/gen_k1_vla_profile_comparison.py
\begin{figure*}[t]
  \centering
  \includegraphics[width=0.95\textwidth]{figures/k1_vla_profile_comparison.pdf}
  \caption{Current-best Merlin warm inference, XNNPACK-kernel swap, and per-operation attribution on the single-core
  SpaceMiT K1. All results use three independent launches with two untimed warmups and five timed
  iterations and pass a full-output golden gate. OpenVLA and RDT2 pass the profiler perturbation
  gate; BitVLA's asterisked mix is diagnostic because its 2.135\% perturbation narrowly
  exceeds the 1.9\% allowance. These are deterministic reduced-depth architecture captures with
  generated parameters, so they characterize compiler/codegen behavior rather than pretrained
  task quality.}
  \label{fig:k1-vla-profile-comparison}
\end{figure*}

% Cross-model single-core INT8 diagnostic. Missing qd8 hybrid cells are explicit.
\begin{figure*}[t]
  \centering
  \includegraphics[width=0.95\textwidth]{figures/k1_model_runtime_comparison.pdf}
  \caption{Single-core warm-inference comparison on the SpaceMiT K1. Absolute latency is shown on
  a log scale at left and normalized to ExecuTorch at right. LSTMNetVIT and TinyLLaMA use the
  current identity-matched diagnostic cells; SmallLLaMA\textsuperscript{\dag} combines the recorded 5.186 ms Merlin
  W8A8 search winner with a 1.881 ms ExecuTorch qd8 warm slope; SmolVLA's hatched Merlin bar is a
  lower bound because it exceeded the 120 s deadline. Red crosses are intentionally not bars: a
  Merlin-runtime plus XNNPACK qd8 backend has not been validated, so no W8A8 hybrid performance is
  claimed. \textsuperscript{\dag}The SmallLLaMA pair is cross-run. SmolVLA is timeout- and
  reference-qualified. These are diagnostic results rather than a paired certification set.}
  \label{fig:k1-model-runtime-comparison}
\end{figure*}

% Single-column (`figure`, \columnwidth) -- every other block here spans two columns.
% Counts are derived from the live corpus, not typed: rerun --derive if the corpus moves.
% Regenerate with: python figures/gen_capsule_generation.py
\begin{figure}[t]
  \centering
  \includegraphics[width=\columnwidth]{figures/capsule_generation.pdf}
  \caption{\textbf{Capability-derived capsule generation.} Merlin derives an executable coverage
  contract before compiler authoring. (a)~Pinned target specifications, RTL-derived capabilities and
  workload requirements define a compiler-independent coverage space: eight axes and 111 required
  cells across four targets. (b)~Representative obligations are materialized as target-bound ISA,
  layer, model-slice and model capsules. (c)~Each capsule exposes the program the compiler compiles
  against, while the numerical golden and the structural instruction-coverage surface remain hidden
  from it. (d)~Per target, the corpus by capsule kind; obligations that could not be materialized
  stay visible as hatched segments rather than shrinking the coverage denominator---four atlas roster
  models whose quantization scheme the host cannot capture, and two host-lane assertions the target
  admits work for and so cannot prove. Public capsules guide Phase~1 authoring; capability-derived
  hidden capsules, counted here but never named, evaluate the frozen compiler afterwards.}
  \label{fig:capsule-generation}
\end{figure}
